\section{Однородные системы уравнений. Теорема о решениях. Теорема о структуре общего решения. Критерий существования нетривиальных решений}

\begin{definition}
Однородной системой линейных уравнений называется система вида $Ax = 0$, то есть:
\[
\begin{cases}
a_{11}x_1 + a_{12}x_2 + \ldots + a_{1n}x_n = 0 \\
a_{21}x_1 + a_{22}x_2 + \ldots + a_{2n}x_n = 0 \\
\vdots \\
a_{m1}x_1 + a_{m2}x_2 + \ldots + a_{mn}x_n = 0
\end{cases}
\]
\end{definition}

\begin{remark}
Однородная система всегда совместна, так как $x = 0$ (нулевой вектор) всегда является решением (тривиальным решением).
\end{remark}

\begin{theorem}[о свойствах решений однородной системы]
Пусть $x^{(1)}$ и $x^{(2)}$ --- решения однородной системы $Ax = 0$. Тогда:
\begin{enumerate}
    \item $x^{(1)} + x^{(2)}$ --- решение;
    \item $\lambda x^{(1)}$ --- решение для любого $\lambda \in \mathbb{R}$.
\end{enumerate}
\end{theorem}

\begin{proof}
\begin{enumerate}
    \item $A(x^{(1)} + x^{(2)}) = Ax^{(1)} + Ax^{(2)} = 0 + 0 = 0$.
    \item $A(\lambda x^{(1)}) = \lambda Ax^{(1)} = \lambda \cdot 0 = 0$.
\end{enumerate}
\hfill$\blacksquare$
\end{proof}

\begin{corollary}
Множество всех решений однородной системы образует линейное подпространство в $\mathbb{R}^n$.
\end{corollary}

\begin{definition}
\textbf{Фундаментальной системой решений} (ФСР) однородной системы называется любой набор линейно независимых решений $y^{(1)}, y^{(2)}, \ldots, y^{(k)}$, через которые линейно выражается любое решение системы.
\end{definition}

\begin{theorem}[о структуре общего решения однородной системы]
Пусть $\operatorname{rang} A = r$, $n$ --- число неизвестных. Тогда:
\begin{enumerate}
    \item Если $r = n$, то система имеет только тривиальное решение $x = 0$.
    \item Если $r < n$, то существует ФСР, состоящая ровно из $n - r$ решений $y^{(1)}, \ldots, y^{(n-r)}$, и общее решение имеет вид:
    \[
    x = c_1 y^{(1)} + c_2 y^{(2)} + \ldots + c_{n-r} y^{(n-r)},
    \]
    где $c_1, \ldots, c_{n-r}$ --- произвольные постоянные.
\end{enumerate}
\end{theorem}

\begin{proof}
\begin{enumerate}
    \item Если $r = n$, то после приведения к ступенчатому виду все переменные --- главные, и из последнего уравнения получаем $x_n = 0$, затем $x_{n-1} = 0$ и так далее.
    
    \item Пусть $r < n$. Приведём систему к ступенчатому виду:
    \[
    \begin{cases}
    a_{11}x_1 + a_{12}x_2 + \ldots + a_{1n}x_n = 0 \\
    a_{22}x_2 + \ldots + a_{2n}x_n = 0 \\
    \vdots \\
    a_{rr}x_r + \ldots + a_{rn}x_n = 0
    \end{cases}
    \]
    Переменные $x_1, \ldots, x_r$ --- главные, $x_{r+1}, \ldots, x_n$ --- свободные. Выразим главные через свободные:
    \[
    \begin{cases}
    x_1 = -c_{1,r+1}x_{r+1} - \ldots - c_{1n}x_n \\
    x_2 = -c_{2,r+1}x_{r+1} - \ldots - c_{2n}x_n \\
    \vdots \\
    x_r = -c_{r,r+1}x_{r+1} - \ldots - c_{rn}x_n
    \end{cases}
    \]
    
    Построим ФСР. Положим последовательно:
    \[
    \begin{array}{c|cccc}
    & x_{r+1} & x_{r+2} & \ldots & x_n \\
    \hline
    y^{(1)} & 1 & 0 & \ldots & 0 \\
    y^{(2)} & 0 & 1 & \ldots & 0 \\
    \vdots & \vdots & \vdots & \ddots & \vdots \\
    y^{(n-r)} & 0 & 0 & \ldots & 1
    \end{array}
    \]
    Для каждого набора свободных переменных находим соответствующие главные. Полученные векторы $y^{(1)}, \ldots, y^{(n-r)}$ линейно независимы (так как их последние $n-r$ координат образуют единичную матрицу) и образуют ФСР.
    
    Любое решение $x$ можно представить как линейную комбинацию $y^{(i)}$ с коэффициентами, равными значениям свободных переменных.
\end{enumerate}
\hfill$\blacksquare$
\end{proof}

\begin{theorem}[критерий существования нетривиальных решений]
Однородная система $Ax = 0$ имеет нетривиальное решение тогда и только тогда, когда:
\begin{enumerate}
    \item $\operatorname{rang} A < n$, или
    \item (для квадратной системы) $\det A = 0$.
\end{enumerate}
\end{theorem}

\begin{proof}
Следует непосредственно из предыдущей теоремы: нетривиальные решения существуют тогда и только тогда, когда есть хотя бы одна свободная переменная, то есть $r < n$. Для квадратной матрицы это эквивалентно $\det A = 0$.
\hfill$\blacksquare$
\end{proof}

\begin{theorem}[о структуре общего решения неоднородной системы]
Пусть система $Ax = b$ совместна, $x^{(\text{ч})}$ --- её частное решение, а $x^{(\text{о})}$ --- общее решение соответствующей однородной системы $Ax = 0$. Тогда общее решение неоднородной системы имеет вид:
\[
x = x^{(\text{о})} + x^{(\text{ч})}.
\]
\end{theorem}

\begin{proof}
\begin{itemize}
    \item[$\Rightarrow$] Пусть $x^*$ --- произвольное решение $Ax = b$. Тогда $A(x^* - x^{(\text{ч})}) = Ax^* - Ax^{(\text{ч})} = b - b = 0$, то есть $x^* - x^{(\text{ч})}$ --- решение однородной системы. Обозначим $x^{(\text{о})} = x^* - x^{(\text{ч})}$, тогда $x^* = x^{(\text{о})} + x^{(\text{ч})}$.
    
    \item[$\Leftarrow$] Пусть $x = x^{(\text{о})} + x^{(\text{ч})}$, где $Ax^{(\text{о})} = 0$. Тогда $Ax = A(x^{(\text{о})} + x^{(\text{ч})}) = Ax^{(\text{о})} + Ax^{(\text{ч})} = 0 + b = b$.
\end{itemize}
\hfill$\blacksquare$
\end{proof}
